\documentclass[blue]{beamer}

\mode<presentation> {
  \setbeamertemplate{background canvas}[vertical shading][bottom=red!10,top=blue!10]

  \usetheme{Boadilla}
  \usefonttheme[onlysmall]{serif}
}


\usepackage{pgf,pgfarrows,pgfnodes,pgfautomata,pgfheaps,pgfshade}
\usepackage{amsmath,amssymb,bm}
\usepackage{colortbl}
\usepackage[portuguese]{babel}
\usepackage[latin1]{inputenc}
\usepackage{enumerate}


\setbeamercovered{dynamic}


\def\R{\mathbb{R}}
\def\Q{\mathbb{Q}}
\def\N{\mathbb{N}}
\newcommand{\ps}[2]{ \langle #1 , #2 \rangle }
\def\1{\mathbf{1}}
\newcommand{\nor}[1]{\left\| #1 \right\|}
\def\leq{\leqslant}
\def\geq{\geqslant}
\def\et{\quad \textrm{e} \quad}
\def\ou{\quad \textrm{ou} \quad}


\DeclareMathOperator\argmax{argmax}
\DeclareMathOperator\Prob{Prob}
\DeclareMathOperator\supp{supp}
\DeclareMathOperator\inte{int}
\DeclareMathOperator\co{co}
\DeclareMathOperator\cl{cl}
\DeclareMathOperator\Dirac{Dirac}

%------------------------------ theoren styles

\theoremstyle{definition}
\newtheorem*{assumption}{Assumption}
\newtheorem*{conjecture}{Conjecture}

\theoremstyle{example}
\newtheorem*{proposition}{Proposition}


%-----------------------------------------

\title[Corporate Finance]{Theory of Corporate Finance}
\author[Victor Filipe]{V. Filipe Martins-da-Rocha}
\institute[FGV EESP]{Escola de Economia de S\~ao Paulo}
\date[April--June, 2026]{Syllabus}
\subject{Economic Theory}


\begin{document}

\frame{\titlepage}




%------------------------------------------------------------------------------------------

\begin{frame}\frametitle{Course Objective}

\begin{itemize}
\item Introducing master and doctoral students to (part of the) theoretical research in corporate finance
\item Emphasis: general equilibrium and contract theory with incomplete information
\item Fundamentals of corporate finance theory
    \begin{itemize}
    \item firm's choice of its capital structure and dividend policy under alternative assumptions
    \item moral hazard and agency problems, adverse selection and signalling
    \end{itemize}
\end{itemize}

\end{frame}

%------------------------------------------------------------------------------------------

\begin{frame}\frametitle{Pre-requisites}

\begin{itemize}
\item Basic notions of general equilibrium and asset pricing
\item Non-cooperative game theory with and without incomplete information
\end{itemize}

\end{frame}

%--------------------------------------------------------------------------------------

\begin{frame}\frametitle{Indicative reading}

The text book for the course is
\begin{itemize}
\item \alert{Tirole, J.: 2006, ``The Theory of Corporate Finance, Princeton University Press'', Princeton and Oxford}
\end{itemize}

Some additional material can be found in
\begin{itemize}
\item Hart, O. D.: 1995, ``Firms Contracts and Finance Structure'', Oxford University Press
\item Smith, C. W.: 1990, ``The Modern Theory of Corporate Finance'', McGraw Hill-Irwin
\end{itemize}

\end{frame}

%--------------------------------------------------------------------------------------

\begin{frame}\frametitle{Other material}

Another Ph.D. level text book is
\begin{itemize}
\item de Matos, J. A.: 2002, ``Theoretical Foundations of Corporate Finance'', Princeton University Press
\end{itemize}

There are many good MBA text books

\begin{itemize}
\item Copeland, T., Weston, J. and Shastri: 2005, ``Financial Theory and Corporate Policy'', Pearson Addison-Wesley
\item Grinblatt, M. and Titman, S.: 1998, ``Financial Markets and Corporate Strategy'', Irwin-McGraw-Hill
\end{itemize}

The standard text book on corporate finance at the MBA level is
\begin{itemize}
\item Brealey, R. A., Myers, S. C. and F., A.: 2008, ``Principles of Corporate Finance'', McGraw Hill-Irwin
\end{itemize}

\end{frame}


%------------------------------------------------------------------------------------------

\begin{frame}\frametitle{Course Outline}

\begin{center}
Part A. Competitive general equilibrium with stock markets
\end{center}

    \begin{enumerate}[1.]
    \item The objective of the firm
    \item The capital structure of the firm
    \item About competitive price expectations
    \item Moral hazard and competitive markets
    \end{enumerate}

\end{frame}


%------------------------------------------------------------------------------------------

\begin{frame}\frametitle{Course Outline}

\begin{center}
Part B. Principal-Agent models
\end{center}


    \begin{enumerate}[1.]
    \item Outside financing: Private benefit and moral hazard
%        \begin{enumerate}
%        \item The role of net worth: a simple model of credit rationing lecture
%        \item Debt overhang lecture
%        \item Borrowing capacity: the equity multiplier lecture
%        \item Nonverifiable income lecture
%        \end{enumerate}
    \item Liquidity and risk management
%        \begin{enumerate}
%        \item The maturity of liabilities lecture
%        \item The liquidity-scale trade-off lecture
%        \end{enumerate}
    \item Corporate financing under asymmetric information
%        \begin{enumerate}
%        \item Implications of the lemons problem and of market breakdown lecture
%        \item Dissipative signals lecture
%        \end{enumerate}
    \item Mergers and acquisitions
    \item Lending relationships and investor activism
    \item Control rights and corporate governance
    \end{enumerate}

\end{frame}

%--------------------------------------------------------------------------------------
\begin{frame}\frametitle{Grading}

Student performance will be evaluated through in-class presentations
\begin{itemize}
  \item presentation of a paper
  \item presentation of a research idea
\end{itemize}

\end{frame}

%--------------------------------------------------------------------------------------
\begin{frame}\frametitle{Personal research interest}

  \centering
  
\includegraphics[width=\textwidth]{figures/magill-quinzii-rochet-ECMA-15.PNG}

\end{frame}

%--------------------------------------------------------------------------------------
\begin{frame}\frametitle{Personal research interest}

  \centering
  \includegraphics[width=\textwidth]{figures/braido-mdr-IER-18.PNG}

\end{frame}


%%--------------------------------------------------------------------------------------
%
%\begin{frame}\frametitle{Professor}
%
%\begin{description}
%
%\item[Nome:] Victor Filipe Martins-da-Rocha
%
%\item[E-mail:] \texttt{victor.rocha@fgv.br}
%
%\item[Sala:] 1117
%
%\item[Tel:] (21) 2559 5871
%
%\item[Web:] \texttt{http://epge.fgv.br/we/VictorRocha}
%
%\end{description}
%
%\end{frame}

\end{document}
